\documentclass[11pt,letterpaper]{article}
\usepackage[T1]{fontenc}
\usepackage[utf8]{inputenc}
\usepackage{lmodern}
\usepackage[margin=1in,headheight=14pt]{geometry}
\usepackage{amsmath,amssymb,amsthm,mathtools,mathrsfs}
\usepackage{microtype,booktabs,tabularx,array,enumitem}
\usepackage[dvipsnames]{xcolor}
\usepackage{fancyhdr,tocloft}
\setlength{\cftbeforesecskip}{4pt}
\usepackage[colorlinks=true,linkcolor=MidnightBlue,citecolor=MidnightBlue,urlcolor=MidnightBlue]{hyperref}
\usepackage[nameinlink,noabbrev]{cleveref}
\hypersetup{pdftitle={Global Divisors and Cohomological Obstructions in Hahn-Coherent Surcomplex Analysis},pdfauthor={Research manuscript},pdfsubject={Surcomplex analysis; Hahn supports; interpolation; Cousin problems; line bundles}}
\pagestyle{fancy}
\fancyhf{}
\fancyhead[L]{\small Hahn-coherent surcomplex analysis}
\fancyhead[R]{\small Global divisors and obstructions}
\fancyfoot[C]{\thepage}
\setlength{\parskip}{0.3em}
\setlength{\parindent}{1.2em}
\setlist{nosep,leftmargin=2em}
\newtheorem{theorem}{Theorem}[section]
\newtheorem{lemma}[theorem]{Lemma}
\newtheorem{proposition}[theorem]{Proposition}
\newtheorem{corollary}[theorem]{Corollary}
\theoremstyle{definition}
\newtheorem{definition}[theorem]{Definition}
\newtheorem{example}[theorem]{Example}
\theoremstyle{remark}
\newtheorem{remark}[theorem]{Remark}
\newcommand{\C}{\mathbb C}
\newcommand{\R}{\mathbb R}
\newcommand{\N}{\mathbb N}
\newcommand{\Z}{\mathbb Z}
\newcommand{\Q}{\mathbb Q}
\newcommand{\No}{\mathbf{No}}
\newcommand{\SC}{\mathbf{SC}}
\newcommand{\OO}{\mathcal O}
\newcommand{\HH}{\mathscr H_\Gamma}
\newcommand{\HP}{\mathscr H_{\Gamma,+}}
\newcommand{\HA}{\mathscr H_{\Gamma,\ge0}}
\newcommand{\mm}{\mathfrak m}
\newcommand{\KG}{K_\Gamma}
\newcommand{\sh}{^{\sharp}}
\newcommand{\supp}{\operatorname{supp}}
\newcommand{\st}{\operatorname{st}}
\newcommand{\ord}{\operatorname{ord}}
\newcommand{\lc}{\operatorname{lc}}
\newcommand{\divr}{\operatorname{div}}
\newcommand{\Piclf}{\operatorname{Pic}_{\mathrm{lf}}}
\newcommand{\Res}{\operatorname{Res}}
\newcommand{\id}{\operatorname{id}}
\newcommand{\dbar}{\bar\partial}
\newcommand{\jet}{\mathcal R}
\newcommand{\BB}{\mathscr B(P)}
\newcommand{\CC}{\mathscr C}
\newcommand{\HCoint}{\mathop{\oint^{\!H}}\nolimits}
\newcommand{\abs}[1]{\left|#1\right|}
\newcommand{\norm}[1]{\left\lVert#1\right\rVert}
\newcolumntype{L}{>{\raggedright\arraybackslash}X}
\numberwithin{equation}{section}
\begin{document}
\hypersetup{pageanchor=false}
\begin{titlepage}
\centering
\vspace*{0.7cm}
{\large\scshape A continuation of the supplied surcomplex framework\par}
\vspace{0.75cm}
{\Huge\bfseries Global Divisors and\par Cohomological Obstructions\par}
\vspace{0.45cm}
{\LARGE in Hahn-Coherent\par Surcomplex Analysis\par}
\vspace{0.75cm}
{\large A sharp realization criterion, deformed interpolation,\par and a line-bundle dichotomy\par}
\vspace{0.75cm}
{\large September 21, 2026\par}
\vspace{0.5cm}
\begin{minipage}{0.93\textwidth}
\begin{abstract}
We study global existence questions for well-ordered Hahn families of ordinary holomorphic functions, evaluated on infinitesimal thickenings of the complex plane. Local preparation alone does not settle these questions: independently permissible local data may require a globally forbidden decreasing sequence of exponents.

We prove an exact criterion for realizing a discrete family of finite infinitesimal zero clusters by one Hahn-coherent entire datum. The criterion is well ordering of the union of supports of the \emph{symmetric coefficients} of the cluster polynomials. Individual root supports need not satisfy it. We then prove a support-restricted Chinese remainder theorem, an all-order deformed Hermite interpolation theorem, and a corresponding Mittag--Leffler theorem for moving poles.

For additive Cousin data on a puncture-and-disk cover, we identify the complete obstruction: the exponents carrying nonremovable ordinary coefficient singularities must form a well-ordered set. This yields explicit injections of $\C^{\N}/\C^{(\N)}$ into first cohomology and into the group of locally free rank-one Hahn-coherent line bundles. That line-bundle group on the ordinary plane vanishes exactly when the exponent group is trivial or ordered-isomorphic to $\Z$. In every other case it contains continuum many linearly independent classes under its natural additive description. The proofs are support-controlled constructions, not convergence arguments in the fine surreal topology.
\end{abstract}
\end{minipage}
\vfill
\begin{minipage}{0.93\textwidth}\small
\textbf{Status.} This manuscript proves new statements in the explicitly specified category inherited from the supplied manuscript. The local preparation machinery is credited to that source, not claimed as new here. No matching global theorem was located in the primary sources consulted, but this is a scoped novelty assessment, not an exhaustive priority certification. The proofs have not been independently refereed or machine checked.
\end{minipage}
\end{titlepage}
\hypersetup{pageanchor=true}
\pagenumbering{roman}
{\small\tableofcontents}
\newpage
\pagenumbering{arabic}

\section{The global problem and the principal results}\label{sec:introduction}

\subsection{Why a global support obstruction is a real analytic issue}
The supplied manuscript \cite{Supplied} constructs a useful middle layer between arbitrary fine-local analytic functions on $\SC=\No[i]$ and the extremely rigid class of functions coherent at every surreal scale. In that middle layer, a function on the infinitesimal thickening of an ordinary domain is represented by
\[
 F(z)=\sum_{\gamma\in S}f_\gamma(z)t^\gamma,
 \qquad f_\gamma\in\OO(U),
 \qquad S\text{ a well-ordered set}.
\]
Here $t^\gamma=\omega^{-\gamma}$ in the surcomplex interpretation. Its coefficients are honest ordinary holomorphic functions on a common ordinary domain. The source proves, among other results, local preparation, exact root counting in a monad, and a moving-pole residue theorem. It also points toward global analytic-set and approximation questions.

The problem addressed here is not to repeat those local results. It is to determine when \emph{infinitely many local prescriptions admit one global coherent datum}. Ordinary complex analysis has very strong answers: arbitrary discrete zero divisors, finite jets, and principal parts can be realized globally on $\C$. The Hahn version has an additional obstruction which those scalar theorems cannot see. The union of all exponents forced by the local conditions may fail to be well ordered.

A first illustration is the proposed list of simple zeros
\begin{equation}\label{eq:impossibleintro}
 z_n=n+t^{1/n},\qquad n=1,2,\ldots .
\end{equation}
Every individual point is legitimate and every finite sublist can be realized. Nevertheless, no Hahn-coherent entire datum has exactly these simple zeros and no others in the thickened plane. In contrast, one can simultaneously realize, at each positive integer $n$, \emph{all} roots of
\begin{equation}\label{eq:possibleintro}
 (z-n)^n-t=0.
\end{equation}
These roots also involve the decreasing exponents $1/n$. The difference is that their symmetric cluster polynomials use only the single exponent $1$. This distinction is the first main theorem, not an artifact of an overly strong sufficient hypothesis.

\subsection{Working category and notation for the main statements}
Fix a set-sized ordered abelian subgroup $\Gamma\subset\No$ and put
\[
 \KG=\C((t^\Gamma)),\qquad
 \HH(U)=\OO(U)((t^\Gamma))
\]
for connected ordinary open sets $U\subset\C$. On disconnected open sets sections are specified componentwise. The resulting object is a sheaf on the \emph{ordinary} plane. Its positive-exponent subsheaf is denoted $\HP$.

For a closed discrete subset $A\subset\C$, prescribe monic polynomials
\begin{equation}\label{eq:clusterintro}
 P_c(u)=u^{m_c}+\sum_{j<m_c}a_{c,j}u^j,
 \qquad a_{c,j}\in\KG\cap\mm,\quad m_c\ge1.
\end{equation}
All roots of $P_c$ in $\SC$ are infinitesimal. They specify the desired roots in the monad of $c$, counted with multiplicity. The \emph{symmetric support} is
\begin{equation}\label{eq:symmetricsupportintro}
 E(P)=\bigcup_{c\in A}\ \bigcup_{j<m_c}\supp(a_{c,j})
 \subset\Gamma_{>0}.
\end{equation}
The family is called admissible when this set is well ordered. There is no upper bound on the multiplicities $m_c$ and no separation assumption on the infinitesimal roots within a cluster.

The results proved below can be summarized as follows.
\begin{enumerate}
\item \textbf{Global divisor realization} (\cref{thm:divisor}). A family \eqref{eq:clusterintro} is the exact zero divisor of a nonzero member of $\HH(\C)$ if and only if it is admissible. A prescribed ordinary leading function with the correct zero divisor can be retained.
\item \textbf{Deformed interpolation and Mittag--Leffler} (\cref{thm:interpolation,thm:ML}). For an admissible cluster family, the correct target of global interpolation is a support-restricted product of the finite algebras $\KG[u]/(P_c)$, not their unrestricted product. The same support condition on proper rational numerators characterizes realizable moving principal parts.
\item \textbf{Cousin obstruction and line bundles} (\cref{thm:cousincriterion,thm:picdichotomy}). On a concrete ordinary cover, additive gluing is equivalent to well ordering of the exponents carrying singular coefficient germs. As a consequence,
\begin{equation}\label{eq:picintro}
 \Piclf(\C,\HH)=0
 \quad\Longleftrightarrow\quad
 \Gamma=\{0\}\ \text{or}\ \Gamma\cong_{\mathrm{ord}}\Z.
\end{equation}
Outside these cases, the left side contains a copy of $\C^{\N}/\C^{(\N)}$ under its additive cohomological description.
\end{enumerate}
Here $\Piclf$ means isomorphism classes of sheaves \emph{locally free of rank one} over $\HH$, with tensor product. This qualification is necessary: the stalks of $\HH$ are not local rings when $\Gamma\ne0$, as we prove in \cref{prop:nonlocal}. We do not silently substitute the larger category of all tensor-invertible modules on an arbitrary ringed space.

\subsection{Provenance and literature boundaries}\label{subsec:provenance}
The definitions of Hahn-coherent evaluation and the local preparation/division strategy come from the supplied manuscript \cite{Supplied}. We reproduce the needed proofs so that the new global arguments can be checked without that file. Normal forms and the ordered-field background are established surreal theory \cite{Gonshor}; the summability algebra is the classical Hahn--Neumann support theory \cite{Neumann}. Algebraic closure after passing to a divisible exponent group follows from the standard Hahn-field theorem \cite[Corollary~4]{Poonen}.

The ordinary interpolation tools are classical, not new results of this article. We give constructive versions on the plane and derive the scalar jet lemma explicitly. McMullen's author-posted notes provide a convenient reference for the ordinary Weierstrass and Mittag--Leffler arguments \cite[\S3]{McMullen}. The cocycle and torsor interpretation of degree-one cohomology is also standard \cite[Tag~02FQ]{Stacks}.

A terminology distinction matters. M\"uller and Strohmaier's \emph{Hahn holomorphic functions} are normally convergent generalized expansions, with a different analytic variable and admissibility condition \cite[Definitions~2.3 and~2.5]{MS}. Our monomials are constant coefficient scales and our ordinary holomorphic functions are their coefficients. No equivalence of these two categories is assumed. Likewise, Ehrlich and Kaplan already study surreal exponential fields and canonical exponentials on suitable surcomplex counterparts \cite{EK}; no global exponential construction is claimed as a contribution here.

The global support criterion, the restricted interpolation algebra, the exact singular-support obstruction, and the ordered-group line-bundle dichotomy are the candidate original contributions. They are not presented as solutions of a previously named published conjecture. The novelty assessment is limited to the sources and searches recorded in the accompanying literature audit. All conclusions concern the stated Hahn-coherent category; they do not settle essential singularities for arbitrary surcomplex extensions or give an all-scale transcendental theory.

\section{Support algebra, evaluation, and the coefficient sheaf}\label{sec:foundations}

\subsection{Hahn summation in a fixed workspace}
A Hahn series over $\Gamma$ is a formal sum
\[
 a=\sum_{\gamma\in S}a_\gamma t^\gamma,
 \qquad a_\gamma\in\C,
\]
whose nonzero support $S$ is well ordered. For $a\ne0$, write $v(a)=\min\supp(a)$. A family of series is \emph{strongly summable} when its total support is well ordered and only finitely many members contribute at each exponent. Addition at an exponent is then finite. All supports, families, and recursions in this article are sets.

For $E\subset\Gamma_{>0}$ define
\[
 E^*=\{e_1+\cdots+e_k:k\ge0,\ e_i\in E\},
\]
where the empty sum is zero.
\begin{lemma}[Support calculus]\label{lem:support}
If $S,T$ are well-ordered subsets of an ordered abelian group, then $S+T$ is well ordered and every exponent has finitely many representations $s+t$. If $E\subset\Gamma_{>0}$ is well ordered, then $E^*$ is well ordered and every exponent is represented by only finitely many finite nondecreasing words in $E$. Consequently, for well-ordered $S$, sums indexed by a term of $S$ and a finite word in $E$ are strongly summable whenever each such word contributes only finitely many terms.
\end{lemma}
\begin{proof}
The first two assertions are the Hahn--Neumann support lemma \cite{Neumann}. For clarity, the last assertion includes two separate requirements. Well ordering follows by applying the first assertion to $S$ and $E^*$. For a fixed output exponent there are finitely many pairs $(s,m)\in S\times E^*$; for each $m$ there are finitely many nondecreasing words with sum $m$, and each finite word has finitely many orderings. Thus there are finitely many contributions before cancellation. This proves exactly strong summability.
\end{proof}
In particular, the phrase ``positive perturbation'' below always refers to a \emph{well-ordered positive support}, not merely to positive valuations of separately chosen terms. In a non-Archimedean ordered group, $n\delta$ need not eventually exceed every positive element. No proof below relies on that false implication.

For a complex vector space $V$ we use $V((t^\Gamma))$ for Hahn families with coefficients in $V$. Every complex-linear map $V\to W$ extends coefficientwise and does not increase support. This observation applies to differentiation, Taylor-jet extraction, ordinary interpolation operators, and local division by a fixed ordinary monomial.

\subsection{Coherent sections and surcomplex evaluation}
For connected nonempty $U$, an element of $\HH(U)$ is
\[
 F=\sum_{\gamma\in S}f_\gamma t^\gamma,
 \qquad f_\gamma\in\OO(U).
\]
The support means the set of exponents with coefficient function not identically zero. Because $\OO(U)$ is a domain, the least exponent $v_H(F)$ satisfies the usual valuation identities for products and sums. Define
\[
 \HA(U)=\{F:v_H(F)\ge0\}\cup\{0\},\qquad
 \HP(U)=\{F:v_H(F)>0\}\cup\{0\}.
\]
The symbol $+$ in $\HP$ refers to exponents, not to the sign of complex coefficients.

Let $\mm\subset\SC$ be the infinitesimals. Every finite surcomplex number has a unique form $c+\epsilon$, where $c\in\C$ and $\epsilon\in\mm$. For an ordinary domain $U$ put
\[
 U\sh=\{c+\epsilon:c\in U,\ \epsilon\in\mm\},
 \qquad \mu(c)=c+\mm.
\]
Define
\begin{equation}\label{eq:evaluation}
 F\sh(c+\epsilon)=
 \sum_{\gamma\in S}\sum_{k\ge0}
 \frac{f_\gamma^{(k)}(c)}{k!}t^\gamma\epsilon^k.
\end{equation}
\begin{proposition}[Evaluation and identity]\label{prop:evaluation}
Formula \eqref{eq:evaluation} is strongly summable and defines an injective algebra homomorphism into functions on $U\sh$. It respects coefficientwise differentiation and formal Taylor expansion in infinitesimal displacements. A coherent datum vanishing on an ordinary nonempty open subset of connected $U$ is zero.
\end{proposition}
\begin{proof}
For $E=\supp(\epsilon)$ the total support lies in $S+E^*$. A power $\epsilon^k$ contributes words of length $k$, so \cref{lem:support} gives both well ordering and finite multiplicity. The ordinary Taylor multiplication identities may therefore be regrouped. This proves the algebra assertion. Evaluation at an ordinary point is $\sum f_\gamma(c)t^\gamma$; uniqueness of normal forms and the ordinary identity theorem give injectivity and the final assertion. Taylor recentering uses the same support calculation with the union of the finitely many displacement supports. The resulting local series has the coefficientwise derivatives. Its fine differentiation follows by factoring the remainder as a square of the increment times a locally bounded Hahn expression, as in the local evaluation argument of \cite{Supplied}.
\end{proof}
The global constructions in this paper are first equalities of coefficient data. We use evaluation only after those data have been assembled. No sequence of partial Hahn sums is asserted to converge in the full fine surreal topology.

\begin{proposition}[The ordinary-base sheaf]\label{prop:sheaf}
Componentwise sections $\HH$, $\HA$, and $\HP$ are sheaves on the ordinary plane. On a connected base, gluing coherent sections produces a single well-ordered support.
\end{proposition}
\begin{proof}
Refine a cover to connected open sets. Equality on an overlap implies coefficientwise equality there. A coefficient function which is nonzero on a connected patch cannot vanish identically on a nonempty open part of that patch. Thus the exact supports on intersecting patches agree. The intersection graph of a connected cover is connected, so one support is shared by all patches. Ordinary holomorphic gluing at each exponent gives the required global section. The nonnegative and positive support conditions are preserved. On disconnected bases the argument applies separately to each component.
\end{proof}
Notice what this does \emph{not} say. Arbitrary local sections need not glue after solving an additive or multiplicative correction problem: the corrections themselves may fail to have a common admissible support. That is the obstruction studied later.

\subsection{Units and positive logarithms}
\begin{proposition}[Unit criterion]\label{prop:units}
On a connected ordinary domain a nonzero $F=t^\lambda(f_\lambda+E)$, with $E\in\HP(U)$, is a unit of $\HH(U)$ if and only if the ordinary function $f_\lambda$ is nowhere zero. In that case
\begin{equation}\label{eq:unitinverse}
 F^{-1}=t^{-\lambda}f_\lambda^{-1}
 \sum_{k\ge0}(-E/f_\lambda)^k.
\end{equation}
Moreover the formal series
\[
 \exp B=\sum_{k\ge0}\frac{B^k}{k!},\qquad
 \log(1+B)=\sum_{k\ge1}\frac{(-1)^{k+1}B^k}{k}
\]
give inverse group isomorphisms $\HP\cong1+\HP$.
\end{proposition}
\begin{proof}
If $FG=1$, the leading coefficient functions multiply to $1$, proving necessity. For sufficiency \cref{lem:support} justifies the geometric series and its multiplication by $F$. The exponential and logarithm have supports in $\supp(B)^*$; all formal identities, including the addition law in a commutative ring, involve finitely many contributions at each exponent and hence evaluate as coefficient identities. They commute with restrictions.
\end{proof}
This logarithm is a logarithm of a \emph{positive-exponent perturbation of one}. It has nothing to do with choosing a global branch of the ordinary logarithm or a global surcomplex exponential at infinite imaginary arguments.

\section{The scalar existence tools, with constructions}\label{sec:scalar}
The following ordinary facts are the analytic input to the global Hahn recursions. Including their proofs makes clear that analytic approximation takes place only at the complex coefficient level.

\subsection{Ordinary divisors and principal parts on the plane}
\begin{lemma}[Scalar Weierstrass and Mittag--Leffler]\label{lem:scalarML}
Let $A\subset\C$ be closed and discrete.
\begin{enumerate}
\item For positive integers $m_c$, there is an entire function $f_0$ whose zeros are precisely the points $c\in A$, with orders $m_c$.
\item For a finite Laurent principal part $p_c$ at every $c\in A$, there is a meromorphic function on $\C$ with exactly those principal parts and no other poles.
\end{enumerate}
\end{lemma}
\begin{proof}
The finite case is immediate. In the infinite case enumerate the nonzero points as $c_n$ with $|c_n|\to\infty$; a possible prescription at zero is treated by one separate factor or summand.

For the first assertion use the canonical factors
\[
 E_p(w)=(1-w)\exp\left(w+\frac{w^2}{2}+\cdots+\frac{w^p}{p}\right).
\]
On $|w|\le1/2$, the branch with value zero at the origin satisfies
\[
 \log E_p(w)=-\sum_{k>p}\frac{w^k}{k},
 \qquad |\log E_p(w)|\le2|w|^{p+1}.
\]
Choose $p_n$ large enough that
\[
 m_{c_n}|\log E_{p_n}(z/c_n)|\le2^{-n}
 \qquad (|z|\le |c_n|/2).
\]
The product of $E_{p_n}(z/c_n)^{m_{c_n}}$ converges normally on compact sets after finitely many factors are separated. Off its specified zeros, the tail is the exponential of a normally convergent sum, hence is nonzero. This proves the asserted divisor.

For the second assertion, $p_{c_n}$ is holomorphic on $|z|<|c_n|$. Choose a polynomial $q_n$ approximating it within $2^{-n}$ on $|z|\le|c_n|/2$, using its Taylor expansion there. Then
\[
 h(z)=\sum_n\bigl(p_{c_n}(z)-q_n(z)\bigr)
\]
converges normally on compact subsets of $\C\setminus A$. Near any one point of $A$ the other terms form a holomorphic function, so $h$ has precisely the prescribed principal part. This is the classical correcting-polynomial construction; see also \cite[Theorem~3.26]{McMullen}.
\end{proof}

\begin{lemma}[Discrete Hermite interpolation]\label{lem:hermite}
For a closed discrete $A\subset\C$ and finite $m_c\ge1$, the map
\begin{equation}\label{eq:scalarjets}
 R_0:\OO(\C)\longrightarrow
 V:=\prod_{c\in A}\C[u]_{<m_c},\qquad
 R_0(f)_c=\sum_{j<m_c}\frac{f^{(j)}(c)}{j!}u^j,
\end{equation}
is onto. It has a complex-linear right inverse $L:V\to\OO(\C)$.
\end{lemma}
\begin{proof}
Choose $f_0$ as in \cref{lem:scalarML}. Given $B_c\in\C[u]_{<m_c}$, take the principal part at $c$ of $B_c(z-c)/f_0(z)$. By the scalar Mittag--Leffler theorem there is meromorphic $h$ with these principal parts. Then $f=f_0h$ is entire and
\[
 f(c+u)-B_c(u)
 =f_0(c+u)\left(h(c+u)-\frac{B_c(u)}{f_0(c+u)}\right)
\]
is divisible by $u^{m_c}$ holomorphically. Thus $R_0f=B$. A surjective linear map of complex vector spaces has a linear section after choosing a basis of its target and preimages of basis vectors.
\end{proof}
The operator $L$ is not claimed to be continuous, unique, or computationally effective on arbitrary infinite input. What matters below is the precise algebraic property $R_0L=\id$. Extending $L$ coefficientwise does not introduce any new Hahn exponent.

\subsection{The scalar additive Cousin theorem}
\begin{lemma}[Scalar additive Cousin theorem on $\C$]\label{lem:scalarCousin}
For any ordinary open cover $(U_i)$ and holomorphic additive cocycle $b_{ij}$, there are $a_i\in\OO(U_i)$ with $b_{ij}=a_i-a_j$. In particular $H^1(\C,\OO)=0$.
\end{lemma}
\begin{proof}
We recall a construction, including the global analytic step. A smooth partition of unity subordinate to the cover, or to a locally finite refinement, gives smooth $v_i$ with $v_i-v_j=b_{ij}$: for the convention $b_{ij}+b_{jk}=b_{ik}$ take $v_i=\sum_k\rho_k b_{ik}$. Their $\dbar v_i$ agree and give a global smooth $(0,1)$-form $a(z)\,d\bar z$.

Every such form on $\C$ is $\dbar u$ for a global smooth $u$. To see this directly, split $a=\sum_{n\ge0}a_n$ by a locally finite smooth annular partition, where each $a_n$ is compactly supported and, for large $n$, vanishes on $|z|<n-1$. Its Cauchy transform
\[
 u_n(z)=\frac1\pi\int_{\C}\frac{a_n(\zeta)}{z-\zeta}\,dA(\zeta)
\]
is smooth, satisfies $\dbar u_n=a_n$, and is holomorphic off the support of $a_n$. These are the ordinary Cauchy--Green identities. For large $n$, approximate the Taylor series of $u_n$ at zero by a polynomial $p_n$ so that all derivatives of order at most $n$ of $u_n-p_n$ are at most $2^{-n}$ on $|z|\le n-2$. For the finitely many small indices no correction is needed. Then
\[
 u=\sum_n(u_n-p_n)
\]
converges in every $C^k$ norm on every compact subset and satisfies $\dbar u=a$. Consequently $a_i=v_i-u$ are holomorphic and have the required differences. If a refinement was used, these expressions still produce a splitting on the original cover, or equivalently the local torsor sections descend. The standard cocycle--torsor interpretation identifies this with $H^1(\C,\OO)=0$ \cite[Tag~02FQ]{Stacks}.
\end{proof}
We shall also use the ordinary fact that every holomorphic line bundle on $\C$ is trivial. One route is the ordinary exponential sheaf sequence
\[
 0\longrightarrow\underline\Z\longrightarrow\OO
 \xrightarrow{\exp(2\pi i\,\cdot)}\OO^\times\longrightarrow1.
\]
The preceding lemma gives $H^1(\C,\OO)=0$, and contractibility of the ordinary plane gives $H^2(\C,\underline\Z)=0$, hence $H^1(\C,\OO^\times)=0$. Similarly $H^1(\C,\underline\Gamma)=0$ for any constant abelian-group sheaf. These are facts about the ordinary topological base, not about the fine topology of $\SC$.

\section{Local preparation and division with uniform support}\label{sec:local}
This section records the local input from \cite{Supplied} in the fixed-$\Gamma$ form needed for infinitely many simultaneous prescriptions. The important extra bookkeeping is that one positive support monoid controls all centers, even when their multiplicities are unbounded.

\subsection{Preparation by an exponent recursion}
For $h\in\OO(D)$ on a disk centered at zero define
\[
 R_mh=\sum_{j<m}\frac{h^{(j)}(0)}{j!}u^j,
 \qquad D_mh=\frac{h-R_mh}{u^m}.
\]
Both maps are complex-linear, preserve holomorphy on the same disk, and satisfy
\begin{equation}\label{eq:ordinarydivision}
 h=R_mh+u^mD_mh.
\end{equation}
They act coefficientwise on Hahn data.

\begin{lemma}[Prepared polynomial and its support]\label{lem:preparation}
Suppose
\[
 H(u)=u^m+E(u),\qquad E\in\HP(D),
 \qquad T=\supp(E),\quad M=T^*.
\]
There is a unique factorization
\begin{equation}\label{eq:localprep}
 H=P Q,\qquad P=u^m+A,\quad \deg_u A<m,\quad Q=1+B,
\end{equation}
where the coefficients of $A$ and the datum $B$ have strictly positive support. The constructed $A,B$ have support in $M_{>0}$, and $Q$ is a coherent unit.
\end{lemma}
\begin{proof}
The equation to solve is $E=A+u^mB+AB$. Proceed by well-founded recursion on $M_{>0}$. At exponent $\nu$, put
\begin{align}
 h_\nu&=E_\nu-
 \sum_{\substack{\alpha+\beta=\nu\\\alpha,\beta>0}}
 A_\alpha B_\beta,\label{eq:preph}\\
 A_\nu&=R_mh_\nu,\qquad B_\nu=D_mh_\nu.\label{eq:prepsteps}
\end{align}
The sum is finite by \cref{lem:support}; its indices are strictly less than $\nu$, so the recursion is legitimate. Equation \eqref{eq:ordinarydivision} verifies the factorization at each exponent. The resulting support is contained in the well-ordered set $M$, and every coefficient remains holomorphic on $D$.

For uniqueness compare two normalized factorizations and take the least exponent of a difference in $A$ or $B$, using the well-ordered union of their supports. At that exponent the quadratic products of differences with positive terms contribute nothing, so
\[
 \Delta A_\nu+u^m\Delta B_\nu=0.
\]
Taylor division forces both differences to vanish, a contradiction. Finally $Q=1+B$ is a unit by \cref{prop:units}.
\end{proof}

\begin{proposition}[Local roots and global uniformity]\label{prop:rootcount}
Let $F\in\HH(U)$ be nonzero, let $\lambda=v_H(F)$, and let $f_\lambda$ be its leading ordinary coefficient. At a zero $c$ of order $m$ of $f_\lambda$, choose a disk with
\[
 q_0(u)=\frac{f_\lambda(c+u)}{u^m}
\]
nowhere zero. Then
\begin{equation}\label{eq:generalprep}
 F(c+u)=t^\lambda q_0(u)P_c(u)Q_c(u),
\end{equation}
where $P_c$ is monic of degree $m$, its lower coefficients are infinitesimal, and $Q_c=1+\text{positive terms}$. The zeros in $\mu(c)$ are exactly the roots of $P_c$, with their polynomial multiplicities. There are no zeros over $c$ if $f_\lambda(c)\ne0$.

If $T=\{\gamma-\lambda:\gamma\in\supp(F),\ \gamma>\lambda\}$, then the lower coefficients of \emph{every} such $P_c$ have support in $T^*_{>0}$.
\end{proposition}
\begin{proof}
Divide $F(c+u)$ by $t^\lambda q_0(u)$ and apply \cref{lem:preparation}. Division by a nonzero ordinary function changes no Hahn exponents, so the same $T^*$ controls every center. The factors $q_0$ and $Q_c$ evaluate to nonzero values throughout the thickened disk.

The polynomial $P_c$ factors over $\SC$, or over a divisible Hahn workspace containing its coefficients \cite[Corollary~4]{Poonen}. Every root $\xi$ is infinitesimal. Otherwise $1/\xi$ is finite and
\[
 \frac{P_c(\xi)}{\xi^m}
 =1+\sum_{j<m}a_{c,j}\xi^{j-m}
\]
is one plus an infinitesimal, not zero. Thus all $m$ roots, counted with multiplicity, lie in $\mm$. Local nonvanishing units do not change zero orders, so \eqref{eq:generalprep} gives the assertion. At a point $c$ with nonzero leading coefficient, evaluation at $c+\epsilon$ retains that nonzero leading value and cannot vanish.
\end{proof}
The normalization identifies $P_c$ intrinsically as the monic polynomial of the entire infinitesimal cluster, not as a choice of labels for its roots.

\subsection{Deformed division and uniform remainders}
\begin{lemma}[Deformed division]\label{lem:division}
Let $P=u^m+A$ with $\deg A<m$ and positive-exponent coefficients. For $N\in\HH(D)$ there are unique $J\in\HH(D)$ and $R\in\KG[u]_{<m}$ such that
\begin{equation}\label{eq:deformeddivision}
 N=PJ+R.
\end{equation}
If $S=\supp(N)$ and $E$ contains the supports of the coefficients of $A$, with $E$ well ordered and positive, then $J$ and $R$ have support in $S+E^*$.
\end{lemma}
\begin{proof}
Set $T(J)=D_m(AJ)$. Applying $D_m$ to \eqref{eq:deformeddivision} requires $(\id+T)J=D_mN$. Define
\begin{equation}\label{eq:divisioninverse}
 J=\sum_{k\ge0}(-T)^kD_mN,
 \qquad R=R_m(N-AJ).
\end{equation}
Every application of $T$ inserts one exponent from $E$, together with a finite Taylor-division operation. The support and finite-contribution assertions therefore follow from \cref{lem:support}. The formal geometric identity gives $(\id+T)J=D_mN$, and \eqref{eq:ordinarydivision} then gives \eqref{eq:deformeddivision}.

For uniqueness, a difference satisfies $P\Delta J+\Delta R=0$. If either difference is nonzero, take the least exponent $\nu$ occurring in the two supports. Multiplication by $A$ cannot contribute at $\nu$, because it raises exponents and all earlier differences vanish. Hence $u^m\Delta J_\nu+\Delta R_\nu=0$, forcing both to be zero by ordinary Taylor division. This is a contradiction.
\end{proof}
Write $\jet_PN=R$ for the remainder. It is $\KG$-linear. Its correction to ordinary Taylor remainder,
\begin{equation}\label{eq:remaindercorrection}
 (\jet_P-R_m)N=-R_m(AJ),
\end{equation}
is supported in $S+E^*_{>0}$, with each term containing at least one exponent from $E$.

For a family $P_c$ whose coefficient supports lie in the \emph{same} $E$, all these conclusions hold with the same controlling support set. Different $m_c$ change the scalar operators but introduce no new exponent. When the coefficients at all centers are viewed as one product vector, the number of exponent decompositions is still finite; it is not necessary to sum over the centers.

\begin{lemma}[Pole-free fractions]\label{lem:polefree}
Let $F,G\in\HH(U)$, $G\ne0$. If the evaluated finite-meromorphic fraction $F/G$ has no actual poles in $U\sh$, then, after filling removable values, it is a member of $\HH(U)$.
\end{lemma}
\begin{proof}
Where the leading coefficient of $G$ is nonzero, use its coherent inverse. Near a leading zero, preparation writes $G$ as a coherent unit times a monic polynomial $P$ with infinitesimal roots. Absorb the unit into $F$ and divide: $N/P=J+R/P$ with $\deg R<\deg P$. Pole-freeness says that $R/P$ has no pole at any root of $P$, hence $P$ divides $R$ over $\SC$. The degree bound gives $R=0$. We have obtained coherent representatives on ordinary neighborhoods of every base point. They agree as fractions on overlaps and hence glue by \cref{prop:sheaf}.
\end{proof}
This argument also proves that two nonzero coherent data with the same actual zero divisor differ by a coherent unit: apply the lemma to both quotients.

\section{An exact global zero-divisor theorem}\label{sec:divisors}

\subsection{Admissible cluster divisors}
Fix a closed discrete $A\subset\C$ and polynomials $P_c$ as in \eqref{eq:clusterintro}. Their zeros define an effective divisor on $\C\sh$: at the finite points $c+\xi$ take the multiplicity of $\xi$ as a root of $P_c$. There are finitely many roots in each monad, and only finitely many relevant monads above a compact ordinary subset. This is the meaning of ``discrete'' throughout this section; fine-topological discreteness alone would be too weak.

\begin{theorem}[Sharp global divisor realization]\label{thm:divisor}
The cluster family $(P_c)_{c\in A}$ is the exact zero divisor of a nonzero $F\in\HH(\C)$ if and only if $E(P)$ in \eqref{eq:symmetricsupportintro} is well ordered.

More precisely, suppose $E=E(P)$ is well ordered and $M=E^*$. For every ordinary entire $f_0$ with zero divisor $\sum_{c\in A}m_c[c]$, there is
\begin{equation}\label{eq:realizer}
 F=f_0+\sum_{\nu\in M_{>0}}f_\nu t^\nu\in\HH(\C)
\end{equation}
with the prescribed cluster divisor. On a fixed ordinary disk around each $c$ it has a factorization
\begin{equation}\label{eq:realizerlocal}
 F(c+u)=P_c(u)Q_c(u),\qquad
 (Q_c)_0(u)=\frac{f_0(c+u)}{u^{m_c}},
\end{equation}
where $Q_c$ is a coherent unit with support in $M$.
\end{theorem}
\begin{proof}
\emph{Necessity.} Let $F$ realize the divisor and normalize by its least monomial $t^\lambda$. Its leading ordinary coefficient has zeros exactly at $A$, of orders $m_c$, by \cref{prop:rootcount}. If $T$ is the positive relative support of $F$, preparation gives normalized cluster polynomials whose lower coefficients lie in $T^*$. Each such polynomial is monic and has exactly the prescribed roots with their multiplicities, so it is the given $P_c$. Thus $E(P)\subset T^*$, and every subset of a well-ordered set is well ordered.

\emph{Sufficiency.} Choose $f_0$ by \cref{lem:scalarML} and choose a disk $D_c$ in the coordinate $u=z-c$ on which $q_{c,0}=f_0(c+u)/u^{m_c}$ is nowhere zero. Write $P_c=u^{m_c}+A_c$. We construct the coefficients of $F$ and
\[
 Q_c=\sum_{\nu\in M}q_{c,\nu}t^\nu
\]
simultaneously by well-founded recursion on $M_{>0}$. Suppose all coefficients below $\nu$ have been chosen. The coefficient equation in \eqref{eq:realizerlocal} is
\begin{equation}\label{eq:globalrecursion}
 f_\nu(c+u)=u^{m_c}q_{c,\nu}(u)+h_{c,\nu}(u),
\end{equation}
where the known function is
\begin{equation}\label{eq:knowncorrection}
 h_{c,\nu}=
 \sum_{\substack{\alpha+\beta=\nu\\\alpha>0,\ \beta\ge0}}
 (A_c)_\alpha\,q_{c,\beta}.
\end{equation}
There are finitely many exponent pairs, and every $\beta$ is less than $\nu$. At every center prescribe the finite ordinary jet
\[
 R_{m_c}\bigl(f_\nu(c+u)\bigr)=R_{m_c}h_{c,\nu}(u).
\]
The discrete Hermite lemma provides one entire function $f_\nu$ satisfying all these conditions, regardless of the number or sizes of the jets. Define
\begin{equation}\label{eq:quotientrecursion}
 q_{c,\nu}(u)=
 \frac{f_\nu(c+u)-h_{c,\nu}(u)}{u^{m_c}}.
\end{equation}
This is holomorphic on the same disk $D_c$, since the numerator has the required vanishing order at zero. No disk shrinking occurs during the recursion.

All coefficients now exist on the single well-ordered support $M$. Equations \eqref{eq:globalrecursion}--\eqref{eq:quotientrecursion} prove the factorizations coefficient by coefficient. The leading coefficient of $Q_c$ is $q_{c,0}$, so it is a unit by \cref{prop:units}. Hence the roots over $c$ are exactly those of $P_c$. There are no other roots because the leading coefficient of $F$ is $f_0$ and \cref{prop:rootcount} excludes all other monads. This proves the exact divisor assertion.
\end{proof}

\begin{corollary}[The ambiguity is exactly a unit]\label{cor:divisortorsor}
All realizers normalized to have the same leading coefficient $f_0$ form a torsor under $1+\HP(\C)=\exp\HP(\C)$. Without that normalization, any two realizers differ by a unique unit
\[
 t^\lambda h_0\exp B,
 \qquad \lambda\in\Gamma,\quad
 h_0\in\OO(\C)^\times,\quad B\in\HP(\C).
\]
\end{corollary}
\begin{proof}
The two quotients are pole free, so both are coherent by \cref{lem:polefree}; they are inverse units. If the leading functions agree, their quotient has leading monomial $1$ and leading coefficient function $1$. Apply \cref{prop:units}. Multiplying any realizer by a unit changes no zero order.
\end{proof}

\subsection{A genuine obstruction despite finite solvability}
\begin{example}[Descending simple displacements]\label{ex:descending}
Take $\Gamma=\Q$, $A=\N_{>0}$, and
\[
 P_n(u)=u-t^{1/n}.
\]
Their symmetric support $\{1/n:n\ge1\}$ is not well ordered. By \cref{thm:divisor}, there is no coherent entire datum with exactly the simple zeros $n+t^{1/n}$. Every finite subfamily is admissible and hence realizable. More generally the prescription $P_n=u-a_nt^{1/n}$, with $a_n\in\C$, is realizable if and only if the sequence $(a_n)$ has finite support. Indeed an infinite subset of these decreasing exponents has no least element.
\end{example}
The obstruction cannot be removed by increasing the exponent workspace: a descending sequence remains descending in every larger ordered group. Thus this is also a nonexistence statement in the full set-supported Hahn-coherent category of the supplied manuscript.

\subsection{Unbounded fractional splitting can nevertheless be realized}
\begin{example}[Symmetric cancellation of forbidden root supports]\label{ex:symmetric}
Take $P_n(u)=u^n-t$ for every positive integer $n$. Its symmetric support is the singleton $\{1\}$. The theorem therefore gives a realizer
\[
 F=f_0+f_1t+f_2t^2+\cdots\in\OO(\C)[[t]].
\]
One explicit permitted leading function is
\begin{equation}\label{eq:leadingproduct}
 f_0(z)=\prod_{n=1}^{\infty}E_2(z/n)^n,
 \qquad E_2(w)=(1-w)\exp(w+w^2/2).
\end{equation}
For $|z|\le R$ and $n>2R$,
\[
 \left|n\log E_2(z/n)\right|
 \le 2n(R/n)^3=2R^3/n^2,
\]
so the ordinary product converges normally and has a zero of order $n$ at each $n$, with no other zeros.

The roots in $\mu(n)$ are
\[
 n+\zeta_n^k t^{1/n},\qquad 0\le k<n,
 \qquad \zeta_n=e^{2\pi i/n}.
\]
Their individual supports have union $\{1/n:n\ge1\}$, which is not well ordered. Nevertheless their complete divisor is realizable with integer Hahn exponents. One may evaluate the resulting datum in $\SC$ or factor all cluster polynomials in a divisible workspace such as $K_\Q$.
\end{example}
There are two different infinite operations in this example. Product \eqref{eq:leadingproduct} is an \emph{ordinary} normally convergent product used to construct one scalar coefficient. Series $F$ is an exact Hahn construction. The proof never takes an infinite Hahn product of the cluster factors over $n$.

\section{A support-restricted Chinese remainder theorem}\label{sec:interpolation}

\subsection{The correct target algebra}
Assume the family $P=(P_c)$ is admissible and set $E=E(P)$, $M=E^*$. Define
\begin{equation}\label{eq:restrictedproduct}
 \BB=\left\{(B_c)_{c\in A}:
 \begin{array}{l}
 B_c\in\KG[u]_{<m_c},\\
 \displaystyle\bigcup_{c\in A}\supp(B_c)\text{ is well ordered}
 \end{array}\right\},
\end{equation}
where $\supp(B_c)$ means the union of the supports of its finitely many polynomial coefficients. Multiplication is performed at each center modulo $P_c$, with the unique degree-$<m_c$ representative.

\begin{proposition}[Restricted product algebra]\label{prop:restrictedalgebra}
Formula \eqref{eq:restrictedproduct} defines a $\KG$-algebra. For families with supports in $S$ and $T$, their product is supported in $S+T+M$. The simultaneous deformed remainder map
\begin{equation}\label{eq:globalremainder}
 \jet_P:\HH(\C)\longrightarrow\BB,
 \qquad F\longmapsto
 \bigl(\jet_{P_c}(F(c+u))\bigr)_{c\in A}
\end{equation}
is a $\KG$-algebra homomorphism.
\end{proposition}
\begin{proof}
Addition has support in the finite union of the two supports. For multiplication apply \cref{lem:division} to $B_cC_c$; its input support lies in $S+T$ and all divisors use the same positive set $E$. The support is therefore contained in $S+T+M$, which is well ordered. Uniqueness of division makes the ordinary quotient-ring identities valid. The same division lemma shows that the image of \eqref{eq:globalremainder} has a common well-ordered support, and uniqueness makes it multiplicative and $\KG$-linear.
\end{proof}
The restriction in \eqref{eq:restrictedproduct} is not merely an optional choice of topology. It is forced by the existence of a single coherent source function.

\subsection{A support-controlled right inverse}
\begin{theorem}[Deformed Hermite interpolation]\label{thm:interpolation}
For an admissible cluster family, prescribed polynomial remainders $B_c\in\KG[u]_{<m_c}$ are simultaneously realized by a member of $\HH(\C)$ if and only if their total coefficient support is well ordered. Thus \eqref{eq:globalremainder} is surjective.

If that support is $S$, a realizing function can be chosen with support in $S+E^*$. More explicitly, let $L$ be a complex-linear right inverse in \cref{lem:hermite}, extended coefficientwise. Put
\[
 N=\jet_P-R_0,
\]
where $R_0$ is also extended coefficientwise. Then a $\KG$-linear right inverse of $\jet_P$ is
\begin{equation}\label{eq:interpolator}
 \mathcal I_P(B)=L\sum_{k\ge0}(-NL)^kB.
\end{equation}
\end{theorem}
\begin{proof}
Necessity follows immediately from uniform deformed division: if $F$ has support $S_F$, all remainders lie in $S_F+E^*$.

For sufficiency identify \eqref{eq:restrictedproduct}, as a vector space, with $V((t^\Gamma))$, where $V$ is the product in \eqref{eq:scalarjets}. A coefficient at one exponent is one element of this product; no sum over all centers is being taken. The scalar map $L$ extends to $V((t^\Gamma))$ without increasing support. It is $\KG$-linear, since multiplication by a fixed Hahn constant is finite convolution at each exponent.

By \eqref{eq:divisioninverse} and \eqref{eq:remaindercorrection}, every term of the operator $N$ inserts a nonempty finite word of exponents from $E$. Thus $NL$ strictly raises the support by such words. In the series in \eqref{eq:interpolator}, a contribution from $(NL)^k$ is indexed by a term of $S$ and a word from $E$ partitioned into $k$ nonempty consecutive blocks. For a fixed output exponent, \cref{lem:support} leaves finitely many such words. Each word has finite length and finitely many block partitions, so only finitely many $k$ can contribute. This proves strong summability, not merely formal plausibility of a Neumann series. Its support is contained in $S+E^*$.

Let $C=\sum_{k\ge0}(-NL)^kB$. The geometric-series identity gives $(\id+NL)C=B$. Since $R_0L=\id$,
\[
 \jet_P(LC)=(R_0+N)LC=(\id+NL)C=B.
\]
This proves existence, the support estimate, and the stated right inverse.
\end{proof}
The proof is valid for arbitrary ordered groups. It never assumes an Archimedean value group, a discrete valuation, a norm estimate for $L$, or a uniform bound on the multiplicities.

\begin{corollary}[Chinese remainder quotient]\label{cor:CRT}
Let $G$ be any realizer of the admissible cluster divisor. Then
\begin{equation}\label{eq:CRT}
 \HH(\C)/(G)\ \cong\ \BB
\end{equation}
as $\KG$-algebras. The quotient map admits a noncanonical $\KG$-linear section.
\end{corollary}
\begin{proof}
Surjectivity and the linear section follow from \cref{thm:interpolation}. If every local remainder of $F$ is zero, then $F$ vanishes to at least the multiplicity of every zero of $G$. Thus $F/G$ has no actual pole and is coherent by \cref{lem:polefree}. Hence the kernel is contained in $(G)$. The reverse inclusion follows from $G=P_cQ_c$ locally. Apply the first isomorphism theorem.
\end{proof}
This is an infinite Chinese remainder statement with a precisely identified image. Replacing $\BB$ by the full product would be false.

\subsection{Bounded raw values need not be interpolable}
\begin{example}[Selecting one root in every growing cluster]\label{ex:selector}
Work in $\Gamma=\Q$ and use the admissible family $P_n(u)=u^n-t$. Write $s_n=t^{1/n}$. At the $n$ roots of $P_n$, prescribe value $1$ at $s_n$ and value $0$ at every $\zeta_n^ks_n$ with $1\le k<n$. The unique polynomial of degree $<n$ with these values is
\begin{equation}\label{eq:selector}
 B_n(u)=\frac1n\sum_{j=0}^{n-1}(u/s_n)^j
 =\frac1n\sum_{j=0}^{n-1}t^{-j/n}u^j.
\end{equation}
Indeed the finite geometric sum is $n$ at the root $1$ of unity and zero at the others. The coefficient supports contain
\[
 -\frac{n-1}{n}=-1+\frac1n,\qquad n\ge2,
\]
an infinite strictly decreasing sequence. Theorem~\ref{thm:interpolation} therefore rules out a global coherent interpolant.

Every prescribed scalar value was either $0$ or $1$, and the underlying zero divisor is realizable even in integer Hahn exponents. The obstruction arises only after converting point values into the intrinsic finite quotient coordinates. Every finite collection of these root-selection conditions is solvable.
\end{example}
If repeated roots occur, values are replaced by finite jets at those roots. The same principle holds: solve the ordinary finite-dimensional polynomial Chinese remainder problem over a splitting field to obtain $B_c$, and then test its global coefficient support. Raw jet coordinates at infinitesimally close points can conceal arbitrarily large negative powers of their separations.

\section{A Mittag--Leffler theorem for moving poles}\label{sec:ML}

\subsection{Principal parts relative to a cluster bound}
Keep an admissible family $P_c$. A prescribed moving principal part at $c$ is a proper rational function
\begin{equation}\label{eq:movingpart}
 \frac{R_c(u)}{P_c(u)},\qquad R_c\in\KG[u]_{<m_c}.
\end{equation}
Two local fractions have the same such principal part when their difference extends to a coherent holomorphic datum on a disk about the ordinary center. By deformed division every local coherent fraction whose poles are bounded by $P_c$ has a unique representative of the form \eqref{eq:movingpart}. This representative packages all individual poles and residues in the cluster, including cancellations.

A \emph{global coherent meromorphic function} here means one fraction $F/G$ with $F,G\in\HH(\C)$, $G\ne0$, interpreted as finite-meromorphic germs with removable values filled. We do not assume that an arbitrary componentwise meromorphic Hahn family belongs to this fraction field.

\begin{theorem}[Sharp moving-pole Mittag--Leffler theorem]\label{thm:ML}
For an admissible pole bound $(P_c)$, the moving principal parts \eqref{eq:movingpart} are realized by a global coherent meromorphic fraction with no poles outside the given bound if and only if
\begin{equation}\label{eq:MLsupport}
 S_R=\bigcup_{c\in A}\supp(R_c)
\end{equation}
is well ordered. Any two solutions differ by a member of $\HH(\C)$.
\end{theorem}
\begin{proof}
Choose the particular divisor realizer $G$ from \cref{thm:divisor}, so that on the disks $D_c$,
\[
 G(c+u)=P_c(u)Q_c(u),
\]
with $Q_c$ and $Q_c^{-1}$ uniformly supported in $E(P)^*$; their leading coefficient functions are ordinary units.

Suppose first that $S_R$ is well ordered. Form the targets
\[
 B_c=\jet_{P_c}(Q_cR_c).
\]
Their total support is contained in $S_R+E(P)^*$ by multiplication and division. By \cref{thm:interpolation} choose $F\in\HH(\C)$ with those local remainders. Then $F-Q_cR_c$ is divisible by $P_c$ coherently, and
\[
 \frac{F(c+u)}{G(c+u)}-\frac{R_c(u)}{P_c(u)}
 =\frac{F(c+u)-Q_c(u)R_c(u)}{P_c(u)Q_c(u)}
\]
is holomorphic on the thickened disk. The quotient $F/G$ has no other poles, because the zeros of $G$ are exactly the specified bound.

Conversely suppose a global fraction $H$ has the required principal parts and no poles outside that bound. Then $GH$ is a pole-free coherent fraction, hence some $F\in\HH(\C)$ by \cref{lem:polefree}. Locally,
\[
 R_c=\jet_{P_c}(F(c+u)/Q_c(u)).
\]
Uniform inversion and deformed division show that all these numerators have support in $\supp(F)+E(P)^*$. Thus $S_R$ is well ordered. Finally the difference of two solutions has no actual poles and is coherent by the same pole-free lemma.
\end{proof}
The pole bound is part of the theorem. Without its admissibility, one cannot assume a single global denominator exists; the divisor theorem gives the exact obstruction to that preliminary step.

\subsection{Residues may have bad global support even when the fraction exists}
\begin{example}[Globally realizable poles with unbounded residue scales]\label{ex:MLresidues}
Again let $P_n(u)=u^n-t$ and prescribe $R_n(u)=1$. Both the divisor support and the numerator support are admissible, so there exists a global coherent meromorphic fraction locally congruent to
\[
 \frac1{(z-n)^n-t}
\]
modulo coherent holomorphic data. At a root $\xi=\zeta_n^ks_n$, its residue is
\begin{equation}\label{eq:individualresidues}
 \frac1{n\xi^{n-1}}
 =\frac{\zeta_n^k}{n}\,t^{-(n-1)/n}.
\end{equation}
The union of these individual residue supports is not well ordered. This does not prevent the global fraction from existing: its proper numerators all equal $1$. For $n>1$, the residues within the $n$th cluster sum to zero, whereas the $n=1$ cluster contributes $1$.
\end{example}
On an ordinary compact contour avoiding leading denominator zeros, the moving-pole residue theorem from \cite{Supplied} therefore applies to only finitely many clusters. It gives the sum of their actual residues. Neither the present theorem nor that contour theorem calls for an infinite sum of residues over the entire plane. The example illustrates why the correct global input consists of proper rational cluster representatives rather than independently enumerated residue coefficients.

\section{The exact additive Cousin obstruction}\label{sec:cousin}
The preceding existence theorems use a common support hypothesis and show that it is sharp. We now identify a gluing obstruction which is invisible to every ordinary coefficientwise existence theorem.

\subsection{A puncture-and-disk cover}
Consider the following ordinary cover of $\C$:
\begin{equation}\label{eq:starcover}
 U_0=\C\setminus\N_{>0},\qquad
 U_n=D(n,1/4)\quad(n\ge1).
\end{equation}
The disks are disjoint. The only nonempty pairwise intersections involving different patches are
\[
 W_n=U_0\cap U_n=U_n\setminus\{n\}.
\]
Thus any family $b_n\in\HH(W_n)$ is additive Cousin data, with the splitting convention
\begin{equation}\label{eq:cousinsplit}
 b_n=a_n-a_0\quad\text{on }W_n.
\end{equation}
Each $b_n$ separately has well-ordered support; their union need not.

Write $b_n=\sum_\gamma b_{n,\gamma}t^\gamma$, interpreting absent coefficients as zero. Define the \emph{singular support}
\begin{equation}\label{eq:singularsupport}
 \Sigma(b)=\{\gamma\in\Gamma:
 \text{some }b_{n,\gamma}\text{ has a nonremovable singularity at }n\}.
\end{equation}
The singularity may be a pole or an ordinary essential singularity. This definition concerns ordinary coefficient functions, not an extension into the infinitesimal monad of a coefficient's essential singularity.

\begin{theorem}[Necessary and sufficient Cousin criterion]\label{thm:cousincriterion}
For the cover \eqref{eq:starcover}, the data $(b_n)$ admit a splitting \eqref{eq:cousinsplit} with $a_0\in\HH(U_0)$ and $a_n\in\HH(U_n)$ if and only if $\Sigma(b)$ is well ordered.

If all $b_n$ have strictly positive support, the same condition is equivalent to a splitting with all $a_i\in\HP(U_i)$.
\end{theorem}
\begin{proof}
\emph{Necessity.} If $b_{n,\gamma}$ has a nonremovable singularity at $n$, then the coefficient $(a_0)_\gamma$ cannot be identically zero on $U_0$: otherwise $b_{n,\gamma}=(a_n)_\gamma$ would extend across $n$. Hence
\[
 \Sigma(b)\subset\supp(a_0),
\]
which is well ordered.

\emph{Sufficiency.} Suppose $\Sigma(b)$ is well ordered. Fix $\gamma\in\Sigma(b)$. The scalar functions $b_{n,\gamma}$ form ordinary additive Cousin data, so the scalar theorem gives $a_{0,\gamma}\in\OO(U_0)$ such that every $b_{n,\gamma}+a_{0,\gamma}$ extends holomorphically across $n$. For $\gamma\notin\Sigma(b)$ set $a_{0,\gamma}=0$; the corresponding $b_{n,\gamma}$ already extend. Let $a_{n,\gamma}$ be these extensions and assemble
\[
 a_0=\sum_{\gamma\in\Sigma(b)}a_{0,\gamma}t^\gamma,
 \qquad a_n=\sum_\gamma a_{n,\gamma}t^\gamma.
\]
The support of $a_0$ is contained in $\Sigma(b)$. For each fixed $n$,
\[
 \supp(a_n)\subset\Sigma(b)\cup\supp(b_n),
\]
a finite union of well-ordered sets. Thus every displayed expression is a legitimate coherent section, and the coefficient equations give \eqref{eq:cousinsplit}. If the input supports are positive, both sets in this last inclusion are positive, so the constructed splitting is positive as well.
\end{proof}

\begin{remark}[An explicit scalar solution on this cover]\label{rem:explicitCousin}
The scalar step can be written without sheaf cohomology. For fixed $\gamma$, take the negative Laurent part $p_n$ of $b_{n,\gamma}$ at $n$. Since its inner Laurent radius is zero, $p_n$ is holomorphic on all of $\C\setminus\{n\}$, even when it is an infinite negative Laurent series. Choose a polynomial $q_n$ with $|p_n-q_n|\le2^{-n}$ on $|z|\le n/2$. Then
\[
 a_{0,\gamma}=-\sum_{n\ge1}(p_n-q_n)
\]
converges normally on compact subsets of $U_0$ and cancels every prescribed negative Laurent part. At each center all other terms extend holomorphically. This is the correcting-polynomial proof of the required scalar Cousin step.
\end{remark}

A union of bad supports is therefore not, by itself, the obstruction. Exponents attached only to removable coefficient germs can stay on the separate disk patches. The exact obstruction is the support that must be carried by the single connected central patch.

\subsection{An explicit infinite-dimensional family of obstructions}
\begin{theorem}[First cohomology detected by descending scales]\label{thm:H1injection}
Suppose $\Gamma$ contains a strictly decreasing sequence
\[
 \gamma_1>\gamma_2>\cdots>0.
\]
Then the cocycles
\begin{equation}\label{eq:obstructioncocycle}
 b_n(a)=\frac{a_n t^{\gamma_n}}{z-n},\qquad a=(a_n)\in\C^{\N_{>0}},
\end{equation}
induce an injective complex-linear map
\begin{equation}\label{eq:cohomologyinjection}
 \C^{\N_{>0}}/\C^{(\N_{>0})}
 \lhook\joinrel\longrightarrow H^1(\C,\HP).
\end{equation}
Here $\C^{(\N_{>0})}$ denotes finitely supported sequences.
\end{theorem}
\begin{proof}
Each component \eqref{eq:obstructioncocycle} is a legitimate positive Hahn section with a one-element support. Its singular support is
\[
 \{\gamma_n:a_n\ne0\}.
\]
An infinite subset of a strictly decreasing sequence indexed by positive integers has no least element. By \cref{thm:cousincriterion}, the cocycle is a coboundary exactly when $(a_n)$ is finitely supported. Linearity is immediate.

To justify the target as sheaf $H^1$, rather than just an unverified cover-dependent quotient, use the standard classification by additive torsors \cite[Tag~02FQ]{Stacks}. A cocycle is zero in sheaf $H^1$ exactly when its glued torsor has a global section. Restricting such a section to the original patches would give a splitting on that original cover. Thus a non-coboundary here cannot disappear on a refinement.
\end{proof}

\begin{corollary}[Size and local invisibility]\label{cor:H1size}
The space in \eqref{eq:cohomologyinjection} has complex dimension $2^{\aleph_0}$. Every individual class constructed there restricts to zero on every bounded ordinary disk, although globally it need not vanish.
\end{corollary}
\begin{proof}
There are only continuum many complex sequences, giving the upper bound on dimension. For each $\lambda\in\C^\times$ consider $(\lambda^n)_{n\ge1}$. Any linear combination of finitely many such sequences with distinct $\lambda$ which is eventually zero has all coefficients zero: apply the Vandermonde determinant to a sufficiently late block of consecutive indices. These classes give continuum many independent vectors in the quotient.

On a bounded disk only finitely many $U_n$ meet it. The corresponding restricted cocycle has finite singular support and splits, either by the theorem or by a finite rational sum. This statement concerns these particular classes. It does not assert vanishing of the sheaf cohomology of every bounded disk for arbitrary data.
\end{proof}

\begin{proposition}[The full additive sheaf is already obstructed in cyclic rank]\label{prop:fulladditive}
For every nonzero $\Gamma$, the group $H^1(\C,\HH)$ contains a copy of $\C^{\N}/\C^{(\N)}$.
\end{proposition}
\begin{proof}
Choose $\delta>0$ in $\Gamma$ and repeat the construction with exponents $\gamma_n=-n\delta$. The general, not necessarily positive, version of \cref{thm:cousincriterion} gives precisely the same kernel.
\end{proof}
This proposition will coexist with a vanishing line-bundle theorem for cyclic $\Gamma$. There is no contradiction: logarithms of units are governed by $\HP$, not by the full Laurent sheaf $\HH$.

\section{Line bundles and the ordered-group dichotomy}\label{sec:linebundles}

\subsection{A necessary distinction between two ringed spaces}
The ordinary-base coefficient sheaf is not the structure sheaf of a locally ringed space in the most naive sense.
\begin{proposition}[Stalks and the integral model]\label{prop:nonlocal}
If $\Gamma\ne0$, the stalk $\mathscr H_{\Gamma,c}$ at an ordinary point $c$ is not local. In contrast, $\mathscr H_{\Gamma,\ge0,c}$ is local, with residue field $\C$ and residue map
\[
 \sum_{\gamma\ge0}f_\gamma t^\gamma\longmapsto f_0(c).
\]
As sheaves of rings, $\HA/\HP\cong\OO$.
\end{proposition}
\begin{proof}
Choose $\delta>0$. In $\mathscr H_{\Gamma,c}$, the germs of $z-c$ and $z-c-t^\delta$ are both nonunits: their leading ordinary coefficient vanishes at $c$, so the unit criterion excludes an inverse on any ordinary neighborhood. Their difference $t^\delta$ is a unit. The nonunits therefore do not form an ideal, which rules out a local ring.

For $\mathscr H_{\Gamma,\ge0,c}$, the displayed map is a surjective ring homomorphism. A germ outside its kernel has exponent-zero coefficient nonzero at $c$, hence that coefficient is nonzero on a smaller disk, and the geometric inverse lies in $\HA$ there. Every germ in the kernel is a nonunit. Thus the kernel is the unique maximal ideal and the residue field is $\C$. Taking coefficient zero is already a surjective sheaf map $\HA\to\OO$ with kernel $\HP$, proving the final assertion.
\end{proof}

\begin{definition}[The line-bundle groups used here]
A Hahn-coherent line bundle over the ordinary plane is a sheaf of $\HH$-modules locally isomorphic to $\HH$ itself. Its isomorphism class lies in $\Piclf(\C,\HH)$. We also use the ordinary Picard group $\Piclf(\C,\HA)$ of locally free rank-one $\HA$-modules. Since $\HA$ has local stalks, this equals its group of tensor-invertible modules and may be denoted simply $\operatorname{Pic}(\C,\HA)$.
\end{definition}
This avoids applying a local-stalk classification to the wrong ringed space. The distinction for general ringed spaces is explicitly relevant to the standard Picard--cohomology statement \cite[Tag~09NT]{Stacks}.

For locally free rank-one modules, transition functions on a trivializing cover take values in the unit sheaf. Changing frames multiplies them by a coboundary. Conversely, unit-valued transition functions glue copies of the structure sheaf. Therefore, for \emph{any} sheaf of commutative rings $\mathscr R$, the category specified by local rank-one freeness gives
\begin{equation}\label{eq:lfclassification}
 \Piclf(\C,\mathscr R)=H^1(\C,\mathscr R^\times).
\end{equation}
This assertion follows directly from the described gluing and torsor construction; it does not assert that all tensor-invertible $\mathscr R$-modules are locally rank-one free when stalks are not local.

\subsection{Splitting the unit sheaves}
\begin{proposition}[Positive logarithms contain all nonclassical line-bundle data]\label{prop:picidentification}
There are isomorphisms of abelian sheaves
\begin{align}
 \HA^\times&\cong\OO^\times\times\HP,\label{eq:Aunits}\\
 \HH^\times&\cong\underline\Gamma\times\OO^\times\times\HP.\label{eq:Hunits}
\end{align}
Consequently extension of scalars from $\HA$ to $\HH$ induces isomorphisms
\begin{equation}\label{eq:picidentification}
 \operatorname{Pic}(\C,\HA)
 \ \cong\ \Piclf(\C,\HH)
 \ \cong\ H^1(\C,\HP).
\end{equation}
\end{proposition}
\begin{proof}
On a connected patch every coherent unit has a unique expression
\[
 F=t^\lambda f_\lambda\exp B,
 \qquad \lambda\in\Gamma,\quad
 f_\lambda\in\OO^\times,\quad B\in\HP.
\]
This is the leading-term decomposition followed by the positive logarithm. Multiplication adds $\lambda$, multiplies $f_\lambda$, and adds $B$, since the rings are commutative. The decomposition respects restrictions. For a unit in $\HA$ one must have $\lambda=0$, giving \eqref{eq:Aunits}; allowing Laurent exponents gives \eqref{eq:Hunits}. The exponent factor is the constant-group sheaf, interpreted componentwise.

Take first cohomology and use \eqref{eq:lfclassification}. The scalar results in \cref{sec:scalar} give
\[
 H^1(\C,\OO^\times)=0,\qquad
 H^1(\C,\underline\Gamma)=0.
\]
Only $H^1(\C,\HP)$ remains, and the inclusion of the two unit decompositions shows that extension of scalars is the stated isomorphism.
\end{proof}
In particular these line-bundle groups have a natural complex-vector-space structure through \eqref{eq:picidentification}. Tensor product corresponds to addition. Every class has an integral model over $\HA$ whose ordinary reduction is a trivial holomorphic line bundle. There is \emph{no} reduction ring map $\HH\to\OO$ sending positive monomials to zero, because those monomials are units in $\HH$; reduction is performed on the $\HA$ model instead.

\subsection{Which exponent groups permit nontrivial bundles?}
\begin{lemma}[A simple ordered-group alternative]\label{lem:groupalternative}
A nonzero ordered abelian group has no infinite strictly decreasing sequence of positive elements if and only if it is ordered-isomorphic to $\Z$.
\end{lemma}
\begin{proof}
The assertion for $\Z$ is immediate. Conversely, if there is no least positive element, choose a smaller positive element repeatedly to construct the forbidden sequence. Thus there is a least positive element $\delta$.

If $\Gamma\ne\Z\delta$, choose $\gamma>0$ not an integer multiple of $\delta$. We claim $\gamma-n\delta>0$ for every integer $n\ge0$. Otherwise choose the first $n$ for which it is nonpositive. Then
\[
 0<\gamma-(n-1)\delta\le\delta.
\]
Minimality of $\delta$ forces equality, so $\gamma=n\delta$, a contradiction. Hence $\gamma,\gamma-\delta,\gamma-2\delta,\ldots$ is a strictly decreasing positive sequence. This also contradicts the hypothesis. Therefore $\Gamma=\Z\delta$.
\end{proof}
A least positive element alone is not enough for the vanishing result. Ordered groups of higher rank may have a least positive element and still contain the descending sequence exhibited in this proof.

\begin{theorem}[Sharp line-bundle dichotomy]\label{thm:picdichotomy}
For every set-sized ordered abelian subgroup $\Gamma\subset\No$,
\begin{equation}\label{eq:picdichotomy}
 \operatorname{Pic}(\C,\HA)=0
 \quad\Longleftrightarrow\quad
 \Piclf(\C,\HH)=0
 \quad\Longleftrightarrow\quad
 \Gamma=\{0\}\ \text{or}\ \Gamma\cong_{\mathrm{ord}}\Z.
\end{equation}
In every other case both groups, identified as in \eqref{eq:picidentification}, contain an embedded complex vector space $\C^{\N}/\C^{(\N)}$ and hence continuum many linearly independent classes. Explicit transition functions realizing the embedding are
\begin{equation}\label{eq:explicitbundles}
 g_{n0}(a)=\exp\left(\frac{a_nt^{\gamma_n}}{z-n}\right)
 \quad\text{on }W_n,
\end{equation}
where $\gamma_1>\gamma_2>\cdots>0$.
\end{theorem}
\begin{proof}
By \cref{prop:picidentification}, it suffices to determine when $H^1(\C,\HP)$ vanishes.

For the trivial group, $\HP=0$. Suppose next $\Gamma=\Z\delta$ with $\delta>0$. On any ordinary open set,
\[
 \HP=\prod_{k\ge1}\OO\,t^{k\delta}
\]
as an additive sheaf: every subset of the positive integer multiples of $\delta$ is well ordered, so there is no further support condition. Given an additive cocycle, split its coefficient at every $k\delta$ by \cref{lem:scalarCousin}. Assemble the resulting local functions. All their supports lie in $\{\delta,2\delta,\ldots\}$, so each assembled function is a legitimate positive coherent section. This splits the original cocycle and proves vanishing. This argument works directly on the original cover; it does not assume that arbitrary infinite products commute with higher sheaf cohomology.

If $\Gamma$ is neither trivial nor ordered-cyclic, \cref{lem:groupalternative} supplies a strictly decreasing positive sequence. The injection of \cref{thm:H1injection} makes $H^1(\C,\HP)$ nonzero and supplies the stated vector subspace. Exponentiation transforms its additive cocycles into \eqref{eq:explicitbundles}; the series are legitimate on each overlap by positive support summability. The unit-sheaf decomposition shows that no extra ordinary unit or monomial change of frame can trivialize these classes unless the original additive cocycle was a coboundary. Finally \cref{cor:H1size} gives the dimension claim.
\end{proof}
For example $\Gamma=\Q$ already exhibits the nontrivial side of the dichotomy. So does every nonzero divisible exponent group. Even a discretely ordered noncyclic group, such as the ordered subgroup $\Z+\Z\omega\subset\No$, lies on that side: $\omega-n$ is a strictly decreasing positive sequence. By contrast $\Gamma=\Z$ has trivial line-bundle group despite the nonzero full additive cohomology in \cref{prop:fulladditive}.

\subsection{Why this does not contradict ordinary complex geometry}
The ordinary plane is contractible and its ordinary holomorphic line bundles are trivial. The nontrivial bundles above are locally trivial for a different coefficient sheaf. Their obstruction is neither an ordinary winding number nor a topological hole in the plane. It is the impossibility of putting infinitely many forced exponents into one well-ordered coherent section on the central patch.

Nor is ``Hahn-coherent'' synonymous with ``coherent $\OO$-module.'' For $\Gamma\ne0$, the sheaf $\HP$ is not even locally finitely generated over $\OO$. At $c$, reduction modulo $z-c$ maps $\mathscr H_{\Gamma,+,c}$ onto the complex vector space of positive Hahn constants. Indeed evaluation of coefficient functions at $c$ gives this map, and if all those values vanish then coefficientwise division by $z-c$ gives its kernel. The quotient contains the infinitely many independent constants $t^{k\delta}$ for any $\delta>0$. A finitely generated $\OO_c$-module would have a finite-dimensional quotient by $z-c$. Thus ordinary theorems for coherent analytic modules cannot be transferred by terminology alone.

\section{Moving divisors as nontrivial line bundles}\label{sec:geometric}
The zero-realization theorem and the cohomology theorem describe the same failure in two complementary languages.

\begin{proposition}[An effective divisor with trivial ordinary reduction but no global equation]\label{prop:Cartierexample}
Let $\gamma_1>\gamma_2>\cdots>0$ in $\Gamma$. There is a locally principal effective divisor on $(\C,\HA)$ whose local surcomplex zeros are the simple points
\[
 n+t^{\gamma_n},\qquad n\ge1,
\]
whose ordinary reduction is a principal ordinary divisor, and whose associated line bundle is nontrivial both over $\HA$ and after extension to $\HH$.
\end{proposition}
\begin{proof}
Choose an ordinary entire $f_0$ with simple zeros exactly at the positive integers. On the cover \eqref{eq:starcover}, take the local equations
\begin{align*}
 h_0(z)&=f_0(z)\quad\text{on }U_0,\\
 h_n(z)&=\frac{f_0(z)}{z-n}\bigl(z-n-t^{\gamma_n}\bigr)
 \quad\text{on }U_n.
\end{align*}
The ordinary quotient $f_0(z)/(z-n)$ is holomorphic and nowhere zero on $U_n$. On $W_n$ the ratio is
\begin{equation}\label{eq:divisortransition}
 g_{n0}=\frac{h_n}{h_0}=1-\frac{t^{\gamma_n}}{z-n},
\end{equation}
a unit of $\HA(W_n)$. The equations are nonzerodivisors and their ratios are units, so they define the asserted effective locally principal divisor and a rank-one line bundle. Coefficient-zero reduction gives the ordinary equation $f_0$ on every patch.

Its positive logarithmic cocycle is
\begin{equation}\label{eq:divisorlog}
 \log g_{n0}=-\sum_{k\ge1}
 \frac{t^{k\gamma_n}}{k(z-n)^k}.
\end{equation}
At exponent $\gamma_n$, the $n$th component has a genuine simple pole. Hence its singular support contains the entire descending sequence $(\gamma_n)$ and is not well ordered. The Cousin criterion and the unit decomposition show that the line bundle is nontrivial in both categories.

Equivalently, a trivializing unit frame $v_i$ with $g_{n0}=v_n/v_0$ would make $h_i/v_i$ glue to a global coherent function with precisely the prescribed divisor. The symmetric support of that divisor is $\{\gamma_n:n\ge1\}$, so \cref{thm:divisor} rules it out.
\end{proof}
This construction has three simultaneous features: all local equations are elementary, the ordinary reduction is globally principal, and the full coherent divisor is not principal. It therefore isolates a genuinely new support obstruction rather than importing a familiar nontrivial ordinary line bundle.

\subsection{What the existence and obstruction results say together}
There are three distinct levels of data in the examples:
\begin{center}
\begin{tabularx}{\textwidth}{@{}LLL@{}}
\toprule
Prescription & Intrinsic global datum & Outcome\\
\midrule
One displaced simple root at each $n$ & $u-t^{1/n}$ & Not realizable: descending symmetric support.\\[0.4em]
All $n$ roots of $u^n-t$ at each $n$ & The single coefficient $-t$ & Realizable, even with integer Hahn exponents.\\[0.4em]
Select one of those roots by a $0/1$ value & Polynomial \eqref{eq:selector} & Not interpolable: descending negative remainder support.\\[0.4em]
Reciprocal principal part at each cluster & Numerator $1$ over $u^n-t$ & Realizable despite descending individual residue scales.\\
\bottomrule
\end{tabularx}
\end{center}
Thus neither ordinary discreteness, finite solvability, bounded raw values, nor the support of individually labeled roots or residues is the correct universal test. The appropriate finite-dimensional invariant must be extracted first: a monic cluster polynomial, a polynomial remainder, or a proper rational numerator. Only then is global support admissibility tested.

\section{Scope, consequences, and further research}\label{sec:scope}

\subsection{What has been resolved in the supplied framework}
The infinite global divisor problem now has an exact answer for the thickened ordinary plane: admissibility of symmetric support is necessary and sufficient. The infinite interpolation image is identified as a concrete restricted product, and moving principal parts have a matching exact criterion. Finally, the Cousin and line-bundle results show that an unrestricted global existence principle is false, with explicit obstructions that survive all coefficientwise scalar solutions.

These conclusions are not consequences of a single ordinary zero-counting theorem. Their proofs couple ordinary global interpolation to transfinite recursion in a common support monoid, while their counterexamples force exponents from infinitely many ordinary locations into one connected patch. The distinction between cyclic and noncyclic exponent groups is exact, not a heuristic contrast between ``small'' and ``large'' infinitesimals.

\subsection{Limits that are part of the statements}
We have worked over the ordinary base $\C$ and its finite surcomplex thickening, not over a hypothetical fine-topological Stein space. In particular no ordinary real-parameter contour is claimed to be fine-continuous, and no transfinite support recursion is treated as a convergent fine-topological iteration. Each positive construction stays inside a set-sized workspace. The nonexistence examples remain nonexistence examples after enlarging that workspace. The explicit nonzero cohomology and line-bundle classes likewise remain nonzero after extension to any larger ordered exponent workspace, since their descending singular supports are unchanged.

The word ``entire'' here means that every coefficient is ordinary entire and the evaluation is defined on $\C\sh$. It does not mean coherence at every surreal scale. The supplied manuscript's all-scale polynomial rigidity is consistent with our examples: an infinite prescribed zero divisor cannot be the divisor of a polynomial, so these realizers cannot satisfy that stronger all-scale hypothesis \cite{Supplied}.

Our line-bundle classification is a vanishing/nonvanishing dichotomy plus a large explicit subspace in the nonvanishing cases. It is not a computation of every element of $H^1(\C,\HP)$ for every noncyclic $\Gamma$. The exact Cousin criterion is for the displayed puncture-and-disk cover; a similarly explicit criterion for arbitrary covers has not been established here. These limitations do not weaken the stated equivalences, whose proofs are complete within their hypotheses.

\subsection{Concrete next questions}
A natural next problem is to characterize the whole group $H^1(\C,\HP)$ in terms of support obstructions, beyond the embedded quotient of sequences. Another is to develop the corresponding finite-rank vector-bundle theory: matrix-valued positive transition functions no longer reduce to an abelian logarithm identity, so noncommutative support recursion must replace the rank-one argument. For global divisor and interpolation results on more general ordinary open Riemann surfaces, the same support recursions should be compared carefully with scalar interpolation and line-bundle obstructions on the base; no unproved generalization is used here.

Several-variable analogues offer a different challenge. The present proof uses finite polynomial remainders in one distinguished local coordinate, and isolated clusters indexed by a discrete ordinary set. Positive-dimensional analytic sets would require coherent parameter-dependent division and compatibility between local finite modules. That is a substantial continuation of the supplied manuscript's multivariable program, rather than an automatic corollary of the one-variable results.

\appendix
\section{Support audit and proof dependencies}\label{app:audit}
For a construction with input support $S$ and admissible positive support $E$, put $M=E^*$. The table records the sets that control the actual infinite operations.
\begin{center}
\begin{tabularx}{\textwidth}{@{}LL@{}}
\toprule
Construction & Controlling support or obstruction\\
\midrule
Evaluation at $c+\epsilon$ & $S+(\supp\epsilon)^*$; words account for Taylor powers.\\[0.3em]
Positive inverse, logarithm, exponential & $E^*$; finitely many words at each exponent.\\[0.3em]
Local preparation & $E^*$, uniformly at all ordinary centers.\\[0.3em]
Global divisor recursion & $E(P)^*$, with one scalar entire coefficient chosen at each exponent.\\[0.3em]
Deformed division & $S+E(P)^*$, independently of cluster multiplicity.\\[0.3em]
Interpolation Neumann series & $S+E(P)^*$; finite words and their finite block partitions.\\[0.3em]
Moving Mittag--Leffler & $S_R+E(P)^*$ after multiplication by local units and division.\\[0.3em]
Cousin splitting & Central support $\Sigma(b)$; disk support contained in $\Sigma(b)\cup\supp b_n$.\\[0.3em]
Nontrivial bundles & A forced infinite descending positive sequence inside $\Sigma(b)$.\\
\bottomrule
\end{tabularx}
\end{center}
Only the ordinary scalar constructions use analytic convergence. They use normal convergence of correcting-polynomial series or smooth convergence of corrected Cauchy transforms. In all Hahn constructions, well ordering \emph{and} finite coefficient multiplicity are checked separately.

The main dependency chain is
\[
 \begin{gathered}
 \text{support calculus + scalar interpolation + local preparation}\\
 \Downarrow\\
 \text{global divisor realization + deformed interpolation}\\
 \Downarrow\\
 \text{moving Mittag--Leffler},
 \end{gathered}
\]
while the cohomological chain is
\[
 \begin{gathered}
 \text{scalar Cousin + exact singular-support criterion}\\
 \Downarrow\\
 \text{descending-scale cohomology classes}\\
 \Downarrow\\
 \text{unit decomposition + ordered-group alternative}\\
 \Downarrow\\
 \text{line-bundle dichotomy}.
 \end{gathered}
\]
The connection between the two chains is the explicit divisor-transition cocycle in \cref{prop:Cartierexample}.

\section{Computational checks and reproducibility}\label{app:checks}
The archive contains \texttt{verify\_examples.py}, which uses exact symbolic arithmetic to check finite instances of the global divisor recursion, local divisibility, root-selector idempotents, and the exponential/logarithm transition identities. It also verifies finite residue identities. The generated \texttt{verification.txt} states the precise tested ranges.

These tests are deliberately separate from the theorem proofs. Finite computations cannot verify well ordering for arbitrary supports, a transfinite recursion, or the universal line-bundle dichotomy. They check signs, coefficients, and finite algebraic instances of the constructions. The article supplies the general mathematical arguments; it does not claim a formal proof-assistant certificate.

The LaTeX file has an internal bibliography and needs no external bibliographic database. Compile it twice with \texttt{pdflatex}, or use \texttt{latexmk -pdf}. The supplementary literature audit records consulted sources and the novelty boundaries, including the distinction between our coefficient sheaf and the normally convergent Hahn-holomorphic theory of M\"uller--Strohmaier.

\begin{thebibliography}{99}
\addcontentsline{toc}{section}{References}

\bibitem{Supplied}
\emph{Surcomplex Analysis: Infinitesimal Calculus, Hahn-Coherent Holomorphy, and Contour Theory}.
Supplied manuscript, dated September 21, 2026, file
\texttt{surcomplex\_analysis(3).tex}.
Source of the Hahn-coherent framework and the local preparation, root-lifting, division, and pole-free-fraction arguments. No publication or authorship attribution is inferred from the filename.

\bibitem{Gonshor}
Harry Gonshor,
\emph{An Introduction to the Theory of Surreal Numbers},
London Mathematical Society Lecture Note Series \textbf{110},
Cambridge University Press, 1986.

\bibitem{Neumann}
B. H. Neumann,
``On ordered division rings,''
\emph{Transactions of the American Mathematical Society}
\textbf{66} (1949), 202--252.

\bibitem{Poonen}
Bjorn Poonen,
``Maximally complete fields,''
\emph{L'Enseignement Math\'ematique} (2) \textbf{39} (1993), 87--106.
\href{https://math.mit.edu/~poonen/papers/amsval.pdf}{Author's full text};
especially Corollary~4 for algebraic closure with divisible value group.

\bibitem{McMullen}
Curtis T. McMullen,
\emph{Advanced Complex Analysis}, Harvard University Math~213a course notes,
November 30, 2023.
\href{https://people.math.harvard.edu/~ctm/home/text/class/harvard/213a/23/html/home/course/course.pdf}{Author's notes}.
See \S3 for ordinary entire and meromorphic functions and Theorem~3.26 for Mittag--Leffler.

\bibitem{MS}
J\"orn M\"uller and Alexander Strohmaier,
``The theory of Hahn meromorphic functions, a holomorphic Fredholm theorem and its applications,''
\emph{Analysis \& PDE} \textbf{7} (2014), no.~3, 745--770.
\href{https://doi.org/10.2140/apde.2014.7.745}{doi:10.2140/apde.2014.7.745};
\href{https://arxiv.org/abs/1205.0236}{arXiv:1205.0236}.

\bibitem{EK}
Philip Ehrlich and Elliot Kaplan,
``Surreal ordered exponential fields,''
\href{https://arxiv.org/abs/2002.07739}{arXiv:2002.07739}.
Used to delimit prior surreal/surcomplex exponential work, not as an input to the global divisor or cohomology proofs.

\bibitem{Stacks}
The Stacks Project Authors,
\emph{The Stacks Project}, Cohomology of Sheaves.
\href{https://stacks.math.columbia.edu/tag/02FQ}{Tag~02FQ: abelian torsors and first cohomology};
\href{https://stacks.math.columbia.edu/tag/09NT}{Tag~09NT: first cohomology and invertible sheaves}.
Consulted September 21, 2026. The local-stalk qualification in the latter is why this article explicitly distinguishes $\Piclf(\C,\HH)$ from the unrestricted tensor-invertible-module terminology.

\end{thebibliography}
\end{document}
